更大的机制诊断面板暴露出固定新进程协议的重复性限制。
30组反馈开关配对中，2组在反馈之前的辅助输出就不同，尽管输入与评价真值哈希一致。
我们保留全部30组原始结果，将这2组标为运行差异混杂；其余28组满足严格上游配对条件。
对这两个输入，各模式再独立运行两次，共增加8个前向任务。
原始结果和重复结果构成12个同模式两两比较，其中6个具有上游差异，
但这两个输入的所有同模式官方AP与召回指标相同。
单复访异常输入的同模式最终logits和mask logits最大绝对差分别达到1.5195和19.3663。
这些差异不能被称为逐张量一致，也不能仅因AP相同就判为无影响。

最早保存的辅助阶段已出现mask logits差异；本轮未定位更早的具体算子，
不将原因归于某一个CUDA内核。
重复实验不替换原始配对，也不将偶然对齐的重复结果挑选为新的干预对照。
因此，全30组反馈开关的汇总差值描述实际运行结果，不能单独证明反馈模块的因果收益。
该发现限制了固定协议下的可重复性主张；新增训练须先解决原生数值一致性问题。
